\begin{abstract}
Data-management systems increasingly change the queries users run in order to meet external requirements: refinement systems rewrite predicates so that the output meets a representation constraint, and consistent query answering drops the tuples involved in integrity violations.
The new output satisfies the requirement, but the user sees only the changed query, not how the population behind the answer has changed.
We study how to explain the difference between two query outputs over sensitive data under differential privacy.
An explanation is a pair of histograms of one attribute over the tuples of the two outputs that fall in a view, a conjunctive region of the data.
We adapt the measures of explanation quality used in prior work (representativeness, coverage, interestingness and diversity) to this setting, bound their sensitivity, and give a two-stage algorithm.
It first privately selects views that are representative of the data and cover the two outputs, then privately selects a diverse set of explanations from these views and releases their histograms with Laplace noise; the whole pipeline is $\epsilon$-differentially private.
On six datasets and $44$ query pairs covering query refinement, consistent query answering and contrasts between groups, the released explanations show how a refinement, a repair or a contrast changes the population, which neither the queries nor the constraints state.
At $\epsilon=1$ they keep $60\%$ of the score of a non-private brute-force search over all candidates, mostly explaining the same attributes, and privacy adds no measurable running time.
\end{abstract}
